% =============================================================================
% Capítulo 7 – Conclusão
% =============================================================================
\chapter{Conclusão}
\label{cap_conclusao}

O presente Trabalho de Conclusão de Curso apresentou o desenvolvimento
do \textbf{SIGEA-GTT} (\textit{Sistema Inteligente de Gestão de Eventos
Adversos – Global Trigger Tool}), um sistema web de software livre voltado
ao ensino e à aplicação da metodologia IHI-GTT no contexto da graduação em
Enfermagem e áreas da Saúde da Universidade Federal de Sergipe (UFS).

\section{Síntese das Contribuições}
\label{sec_contribuicoes}

O SIGEA-GTT entregou um conjunto coeso e interdisciplinar de contribuições
nos planos técnico, pedagógico e científico:

\begin{enumerate}
    \item \textbf{Contribuição Técnica}: A plataforma web foi construída sobre
    uma arquitetura moderna, robusta e de alto desempenho — Spring Boot 3.3
    (Java 21 LTS), Angular 18 com \textit{Standalone Components} e \textit{Signals},
    PostgreSQL 16 e Flyway — atendendo a todos os 25 Requisitos Funcionais (RF01–RF25)
    e 10 Requisitos Não Funcionais (RNF01–RNF10) especificados. Atingiu o
    \textit{Quality Gate} de 100\% de cobertura de testes automatizados
    (JUnit 5 + Jacoco no backend; Jest no frontend), 100\% de geração de
    documentação automática (Javadoc, SpringDoc OpenAPI 3 e Compodoc) e
    conformidade plena com as normas LGPD e de bioética em pesquisa;

    \item \textbf{Contribuição Pedagógica}: O módulo educacional implementado
    oferece ao docente de saúde um ambiente totalmente digital para criar
    prontuários clínicos simulados de alta fidelidade, conduzir turmas
    acadêmicas, gerenciar atividades práticas com cronômetro estrito de 20
    minutos e avaliar formativamente cada estudante com nota e parecer
    pedagógico estruturado. Para o discente, o sistema proporciona uma
    experiência segura e imersiva de auditoria de prontuários fictícios,
    eliminando o risco de danos a pacientes reais e superando as limitações
    impostas pela LGPD no acesso a prontuários hospitalares autênticos;

    \item \textbf{Contribuição Científica}: A integração inédita de seis
    Ferramentas da Qualidade (\textit{Ishikawa}, \textit{GUT}, \textit{5W2H},
    \textit{PDCA/PDSA}, \textit{SWOT} e \textit{Brainstorming}) em uma
    plataforma de ensino em saúde, ancorada no protocolo internacional
    IHI-GTT, representa um avanço metodológico na Informática em Saúde
    Educacional. O sistema gera automaticamente os indicadores epidemiológicos
    do IHI ($TI_{1000}$, $TI_{100}$, $P_{EA}$), viabilizando pesquisas
    quantitativas sobre a incidência de eventos adversos em cenários simulados.
\end{enumerate}

\section{Avaliação do Cumprimento dos Objetivos}
\label{sec_avaliacao_objetivos}

Os objetivos específicos declarados na \autoref{cap_introducao} foram integralmente
alcançados ao longo das quatro fases de desenvolvimento:

\begin{itemize}
    \item \checkmark\ Implementação do controle de acesso RBAC com três perfis
    (Admin, Professor, Estudante) e autenticação JWT \textit{stateless};
    \item \checkmark\ Catálogo completo dos 53 Gatilhos Clínicos GTT agrupados
    nos sete Módulos Assistenciais e das 9 Categorias de Gravidade NCC MERP;
    \item \checkmark\ Prontuário Clínico Simulado persistido em JSONB com dados
    multiprofissionais, prescrições, exames e evolução;
    \item \checkmark\ Cronômetro de 20 minutos para controle estrito do tempo de
    revisão do prontuário, conforme protocolo IHI;
    \item \checkmark\ Seis Ferramentas da Qualidade interativas implementadas no
    frontend Angular com validação reativa em tempo real;
    \item \checkmark\ Dashboard analítico com cálculo automático dos indicadores
    IHI e métricas pedagógicas da turma;
    \item \checkmark\ Interface responsiva, ergonômica e com paleta de cores
    hospitalares validada em todas as quatro resoluções homologadas.
\end{itemize}

\section{Limitações Identificadas}
\label{sec_limitacoes}

Apesar dos resultados alcançados, algumas limitações foram identificadas
durante o desenvolvimento:

\begin{itemize}
    \item O sistema não realiza integração direta com sistemas de
    informação hospitalar reais (HIS/LIS), uma decisão intencional de
    conformidade com a LGPD e com a ética em pesquisa educacional.
    A inclusão de casos clínicos reais anonimizados é um vetor de
    expansão para trabalhos futuros;
    \item O módulo de geração automática de casos clínicos por
    Inteligência Artificial (IA Generativa) foi identificado como uma
    oportunidade, mas não foi implementado nesta versão por exceder o
    escopo do TCC;
    \item A avaliação empírica com usuários reais (discentes e docentes de
    Enfermagem da UFS) não foi realizada neste trabalho, limitando-se à
    validação técnica pelo orientador e pela coorientadora.
\end{itemize}

\section{Trabalhos Futuros}
\label{sec_trabalhos_futuros}

Com base nas limitações identificadas e nos desdobramentos naturais
desta pesquisa, propõem-se os seguintes encaminhamentos futuros:

\begin{enumerate}
    \item \textbf{Validação Empírica com Usuários Reais}: Conduzir um
    estudo experimental controlado com turmas reais do Departamento de
    Enfermagem da UFS, mensurando ganhos de aprendizagem, satisfação de
    usabilidade (SUS) e acurácia na identificação de gatilhos antes e
    após o uso do SIGEA-GTT;
    \item \textbf{Integração com IA Generativa}: Incorporar um agente de
    IA (LLM) para geração automatizada de prontuários clínicos fictícios
    plausíveis, ampliando o banco de casos sem sobrecarga para o docente;
    \item \textbf{Exportação e Interoperabilidade HL7/FHIR}: Desenvolver
    adaptadores de interoperabilidade com padrões internacionais de saúde
    digital (HL7 FHIR R4), viabilizando a importação de casos clínicos
    anonimizados de hospitais conveniados à UFS;
    \item \textbf{Aplicativo Móvel}: Embora o sistema não seja
    \textit{mobile-first}, o desenvolvimento de uma versão progressiva
    (PWA) para revisão de prontuários em tablets de beira de leito é uma
    oportunidade relevante para o contexto hospitalar universitário;
    \item \textbf{Publicação Científica}: Submissão dos resultados a
    periódicos indexados da área de Informática em Saúde (ex: JMIR,
    \textit{International Journal of Medical Informatics}) e ao Congresso
    Brasileiro de Informática em Saúde (CBIS).
\end{enumerate}

\section{Considerações Finais}
\label{sec_consideracoes_finais}

O SIGEA-GTT constitui uma contribuição concreta e de impacto social para
a formação dos futuros profissionais de saúde da UFS. Ao unir rigor
científico da Engenharia de Computação com as demandas reais da educação
em Enfermagem, o sistema demonstra que a interdisciplinaridade — quando
orientada por metodologias sólidas e princípios éticos claros — produz
resultados superiores à soma de suas partes.

A frase que norteia o trabalho permanece: \textit{``Errar em ambiente
simulado é aprender; errar no paciente é imperdoável''}. O SIGEA-GTT
existe para que o segundo nunca ocorra por falta do primeiro.
